\section{User interface}
\label{sec:user-interface}
\subsection{Runtime configuration}
\label{ssec:runtime-configuration}

\begin{docCommand}[doc new=2025-11-19]{liquidsetup}{\marg{options}}
  \index{setup}
  Applies changes to \pkg{liquidtex} configuration at runtime. %
  The options tagged as \enquote{preamble only} (\preambleonly) in
  \cref{sec:package-options} can only be used before the start of the
  document. %
\end{docCommand}

\begin{dbexample}{Setting configuration values in preamble}{}
  \begin{preamblecode}
    \liquidsetup{
      outdir = .liquid-cache,
      use-auxdir = false,
      verbose = true,
      data = info.yml
    }
  \end{preamblecode}
\end{dbexample}

\subsection{Liquid template rendering}

\begin{docCommand}[doc new=2025-11-19]{inputliquid}{\oarg{data
      file}\marg{template file}}
  This command will render a Liquid template and input the result. %
  \begin{itemize}
  \item \meta{template file}: mandatory path to a Liquid template;
  \item \meta{data file}: optional override of the global \refKey{data} setting.
  \end{itemize}
\end{docCommand}

\begin{dbexample}{Rendering templates with an external engine}{}
\begin{preamblecode}
  \usepackage[
    engine      = external,
    engine-conf = { path = liquidtex-cli },
    data        = cv.yml,
    outdir      = .liquid-cache
    ]{liquidtex}
\end{preamblecode}
\begin{bodycode}
  % Uses cv.yml
  \inputliquid{templates/header.liq}

  % Override data for this template
  \inputliquid[teaching.yml]{templates/teaching-entry.liq}
\end{bodycode}
\end{dbexample}
%%% Local Variables:
%%% mode: LaTeX
%%% TeX-master: "../liquidtex"
%%% End:
